\chapter{Layers and Literal Depth}
\label{ch:layers}

\section{One Span per Column}

A sketch layer is a slab, and the one sentence that governs it is this:

\begin{quote}
A layer holds exactly one span per column, and where two adds contest a cell
\textbf{the taller replaces the shorter outright, floor included}.
\end{quote}

Those last five words are, as the repository's own account of the technique puts
it, the whole of what makes sculpture possible. They mean that a stack of nested
shapes writes an arbitrary height field with no subtraction anywhere and no
per-column authoring. Draw a wide low disc, then a narrower taller one on top of
it, then a narrower taller one again, and the result is a stepped mound. Nothing
computed a mound; at every column the tallest add present is simply the one that
survives.

It also means the word \emph{slab} is misleading:

\begin{quote}
A layer is not a flat slab. It is one arbitrary \emph{height field}, a
\id{(floor, top)} pair per column, and a stack of layers is a set of those fields
with air between them. Any solid whatever can be written in that form.
\end{quote}

\subsection{The Algebra of a Single Layer}

Within a layer, a group of shapes resolves as
\[
\bigl((\text{adds} - \text{subtracts}) \cup \text{override-adds}\bigr) - \text{override-subtracts}.
\]

Two consequences follow, and both catch authors out. An \emph{ordinary} add loses
to a subtract over the same column regardless of which was authored first, because
the subtract side of the algebra runs before either shape's override flag is
looked at. And the two add sets each settle among themselves by height before the
override plane overwrites the ordinary one, so an override add replaces the column
it lands on whatever its height is. That is what puts a one-block floor
inside a twelve-block wall, and a threshold through a doorway, with neither shape
cut against the other.

\begin{ruling}{The Limit of Nesting}
A bowl falls inward, and the outer disc that should keep only its own ring is also
the tallest thing over the middle, so nested discs come out as a flat plate. A
falling field needs shapes that do not overlap at all.
\end{ruling}

\section{A Storey as Rock with Holes}

This is the part that most often has to be unlearned. The instinct, drawing a
second storey under a map, is to draw the rooms. That produces nothing, and the
reason is the one-span-per-column rule:

\begin{quote}
A sketch layer carries one span per column, so a room drawn on its own leaves every
other column of that storey empty. The board's lower half is air and the ground
above it floats. What is stated is the \textbf{rock}, over every column the board
has, and the rooms are the holes in it.
\end{quote}

Building an undercroft therefore takes three coordinated moves.

\subsection{The Lift}

The ground has to be raised to make room underneath it. That means the plan states
a higher surface, \id{"globals": \{ "surface": 22 \}}, and the finish thins the
slab to its top courses, leaving the rest free.

Moving only one of the two fails in a way that is confusing to diagnose. If the
surface is lifted without telling the plan, the spawns stay at their old height
while the ground goes up, and the buildability check then reports every placement
as standing over open void, because the marker is a dozen courses under the ground
it names.

\subsection{The Layer Below}

\begin{lstlisting}[language=json]
{ "id": "under", "name": "Undercroft", "base_y": 0, "below": true,
  "shapes": [
    { "id": "rk1",  "type": "rectangle", "operation": "add",
      "floor": 0, "base_height": 14, "theme": "rock",
      "min_x": -50, "min_z": -50, "max_x":  50, "max_z": -30 },
    { "id": "rk2",  "...": "min_x -50, min_z -30, max_x -30, max_z 30" },
    { "id": "rk3",  "...": "min_x  30, min_z -30, max_x  50, max_z 30" },
    { "id": "rk4",  "...": "min_x -50, min_z  30, max_x  50, max_z 50" },
    { "id": "hall", "type": "rectangle", "operation": "add",
      "floor": 0, "base_height": 8, "theme": "room",
      "min_x": -30, "min_z": -30, "max_x": 30, "max_z": 30 }
  ],
  "groups": [ { "id": "under", "mirrors": false, "shapeIds": [ "..." ] } ] }
\end{lstlisting}

\id{below: true} inserts the layer under the compiled ground rather than over it.
Every shape in it is an \textbf{add}: the rock is banded round its hole rather
than cut with a subtract. Four rectangles ring the hall; the hall is a fifth.

The depth arithmetic is then by construction rather than by any field. The rock
runs from floor 0 for fourteen courses, meeting the landmass's underside at
$y=14$. The hall is a shorter span in the same columns, floor 0 and eight courses,
so its floor is stood on at $y=7$ and it has six courses of headroom under the
rock's ceiling. Stated generally: a storey's headroom is the gap between two spans
and needs no field at all.

\subsection{Stamping Without Mirroring}

The group carries \id{mirrors: false}, and that is not a detail.

\begin{quote}
The storey covers the whole board, so its own \id{rot\_180} image lies over it;
fanning it would put every column's ground back twice. A layer that covers only one
half is fanned and one that covers both is stamped once.
\end{quote}

The cost of stamping once is that the layer then has to be its own image, which is
to say symmetric by construction rather than by machinery.

\section{Three Failure Modes}

\subsection{Subtracts Across Layers}

The natural way to make a hall is to cut it out of a slab. The studio refuses
that, twice:

\begin{lstlisting}
POST /map/from-documents   422
SK13  's0' fills 3600 column(s) that 'cut1' takes away -- from (-30, -30)
      -- 's0' is on layer 'ground' and the subtract on 'under', and a
      subtract reaches only the layer it is on
SK13  'descent' fills 56 column(s) that 'cut1' takes away -- from (-22, 2)
\end{lstlisting}

The reason is semantic rather than technical. A subtract states that the column is
\emph{empty}, and the landmass over it then fills what the storey below said was
void. So the recipe is banding: split into bands at every hole edge along one
axis, split each band along the other, and drop the parts that fall inside a hole.

The sculptor's version of the same move avoids the problem at the outline level:
an outline is filled even--odd, so a ring is drawn as the outer ellipse, slit
inward, the inner ellipse the other way round, and closed: one polygon, no
subtract, no refusal.

\subsection{Silent Occlusion}

This is the failure to be most afraid of, because everything reports success:

\begin{quote}
Writing the rock as one rectangle over the whole board and the hall as a shorter
add in the same columns compiles, builds and exports with the gate \textbf{open},
and there is no hall.
\end{quote}

No gate fires, and the studio is right not to fire one: two adds at one floor are
ordinary ground, and the taller winning is exactly what \emph{height is the tallest
add} means. That is precisely why the banding is the technique rather than a detail
of it.

\subsection{Verification by Count}

A missing column in a storey is invisible in every render, because the rock above
and below it hides the gap. So it is not inspected; it is counted. The undercroft
showcase states its check as a number: over the board's 10,000 columns, none is
bare and none carries two spans.

The honest cost of the technique shows up when the holes get complicated. A hall is
four bands. A maze of a hundred corridors and links leaves \textbf{213 rock bands}
between them.

\begin{ruling}{Banding Versus Overlay}
The rock is banded round the hall because the two differ in \emph{height},
fourteen courses against eight, and on one layer the taller add wins the column
outright. A basin and a hall that are the same eight courses in the same columns
win nothing from each other: the geometry is identical either way and only the
paint differs. Getting it the wrong way round is silent in both directions: a
banded-out basin is the same world, and a shorter add inside a taller one is no
room at all.
\end{ruling}

\section{Override and Floor}

An underpass makes the mechanism measurable. A deck and a cut occupy the same
cells. Without \id{override}, a column read at the deck answers \emph{0 solid
blocks, void}: the ordinary add lost to the subtract.

With \id{override: true} but no \id{floor}, the studio refuses:

\begin{lstlisting}
POST /map/from-documents   422
[refusal] SK13  'deck' fills 260 column(s) that 'gullet-cut' takes away
                -- from (-28, -5)
\end{lstlisting}

The reason is that \id{floor} defaults to zero, the very bottom of the shape's
column, so an override add with no stated floor spans from the world's floor to
its own top, and once again fills what the subtract declared void.

\id{floor} is the field that decides, and it decides by setting where the column's
own span \emph{starts}, not by any relation to the ground it replaces. With
\id{"floor": 13, "base\_height": 4} the same column reads four blocks, $y13$ to
$y16$, and nothing recorded beneath. The rule that separates a lid from a fill is
then exact: an override add resting \emph{above} the subtract's floor moves that
span up and records nothing beneath it, and the gate is silent on it; the same
shape at the subtract's own floor is the refusal above. The difference between the
two documents is one field.

\subsection{The Scope of the Constraint}

It is often said that the studio has no floor underneath a shape. That is not
quite it, and the accurate version is more useful:

\begin{quote}
A column carries one span \emph{per sketch layer}, not one span. \ldots So the
honest statement is narrower and more useful: \textbf{one layer cannot stack, and
stacking is what layers are for.}
\end{quote}

\subsection{The Stairwell}

The two mechanisms combine in the move that connects a storey to the one above it:

\begin{lstlisting}[language=json]
{ "id": "descent", "type": "polygon", "operation": "add",
  "override": true, "keepClear": true,
  "floor": 8, "base_height": 14, "theme": "stair",
  "height_mode": "level", "skirt": 0,
  "vertices":       [[-22,2],[-18,2],[-18,16],[-22,16]],
  "anchor_heights": [   1,     1,      14,      14   ] }
\end{lstlisting}

Two numbers matter and they are the same number. The flight's \id{floor} is 8,
the hall floor's own top, so it rests on the storey below rather than driving
through it. And an override add is what cuts the shaft; a subtract there would be
the refusal above again.

One rasterizer fact is worth carrying away from this, because it explains a class
of ugly output that looks like a bug. A ramp at one course per cell only rasterizes
cleanly because the surface is sampled at the cell's centre and \emph{floored} into
the column. Rounding to nearest instead puts every sample of a 1:1 grade exactly on
the boundary, and the flight comes out as a beat of noughts and twos.

\section{Paint Order in a Stack}

Layers are written bottom-up and the painter walks them in document order, which
produces one failure worth stating plainly: a storey listed after one that stands
over it finds no stone left to paint. An undercroft goes \emph{before} the compiled
ground in the document, and \id{below} is what puts it there.

Two further properties of a stacked board catch authors out. A plain layer's paint
bands run from the bedrock course whatever its \id{base\_y}, since only a layer
marked as a \emph{made} thing is painted over its own span, so every ground theme
on such a board wants a fill that is not trivial. And a placement names its storey
explicitly; naming none takes the top surface, which on a roofed goal is the roof.

\section{Made Layers}

A layer marked \id{kind: "made"} is out of the stacking rules: it is painted over
its own span, and it is exempted from the walk that checks a board's pairs and the
walk that checks its reachability. The reason is stated as a fact about the shape
rather than as a convenience. A solid standing in a hill has no gap to lose, and
its roof is not a stair somebody forgot.

A made layer also carries \id{part\_of}, naming the thing a run of layers is sliced
from, and may carry \id{seat: "ground"}, which settles the run onto the terrain
together rather than each slice finding its own height. That is how a sculpted
landmark, whether a dome, a spire, a tapered tower or a gatehouse, is authored: as
ordinary sketch shapes, circles and polygons with a floor and a height, rather than
as stamped block soup.
